\subsection{Multi-exit Networks and Early Exit Policy} % 改叫early exit networks?.
The multi-exit network improves accuracy at the shallow layer by adding loss constraints to internal exit branches. Early exit branches reduce computational resources while enhancing accuracy \cite{kaya2019shallow}. In the inference phase, instances stop inference and leave the model from different branches based on pre-defined rules. Different models use various exit strategies, such as entropy-based or confidence-based scoring, to exit the model\cite{teerapittayanon2016branchynet,kaya2019shallow}. Recently, power\cite{shrestha2022occamnets} shows that multi-exit structure can mitigate dataset bias and improve model predicted. In this paper, we uses the confidence-based early exit strategy to ensure that the early exit instance is confident enough to be correct. To our best knowledge, we are the first work that leverages the early exit technique to improve fairness.

% The multi-exit network is designed by putting additional loss constraints at internal exit branches (internal classifier) to increase the accuracy at shallow layer. With the early exit scheme, early exit branches reduce the computational resource during the model inference time while enhancing the model's accuracy. In the inference phase, the instances will stop inferencing and leave the model from different branches following the pre-defined early exit rules or criteria. BranchyNet \cite{teerapittayanon2016branchynet} calculates entropy at each branch after obtaining the result and exits if the entropy of the predicted result is less than the threshold value. Shallow-Deep Networks \cite{kaya2019shallow} calculates the softmax score of each internal classifier's prediction and takes the maximum probability value as the confidence score. Once the score exceeds the threshold during the forward passing, the instance will exit from the branch prematurely.This model further mitigates the ``overthinking'' problem of deep neural networks. \cite{zhou2020bert} makes predictions using adjacent layers and stops inference when the predicted value of the internal classifier remains constant in a given inference unit times. \cite{shrestha2022occamnets} leverage the multi-exit structure to mitigate the dataset bias and improving model predicted power.

% In this paper, we borrow the confidence-based early exit strategy proposed by \cite{kaya2019shallow}to make sure the the early exit instance is confident enough to be correct. To our best knowledge, we are the first work that leverages the early exit technique to improve fairness.


% \subsection{Bias Mitigation Methods}
Bias mitigation methods aim to reduce native bias in datasets and avoid unfair predictions. There are three main approaches to debiasing: pre-processing, in-processing, and post-processing.

Pre-processing strategies adjust training data before training \cite{ngxande2020bias,lu2020gender} or assign different weights to different data samples to suppress the sensitive information during training \cite{kamiran2012data}, but may not be suitable for dermatological data, where the sensitive information needed for diagnosis (such as diseased areas) is present in the images.

In-processing modifies the model architecture, training strategy, and loss function to remove sensitive information\cite{wang2020towards} and achieve fairness, such as adversarial training \cite{alvi2018turning,zhang2018mitigating} or regularization-based \cite{quadrianto2019discovering,jung2021fair} methods. Recently, pruning \cite{wu2022fairprune} techniques have also been used to achieve fairness in dermatological disease diagnosis. However, implementing these methods may result in decreased accuracy for both groups and does not guarantee explicit protection for the unprivileged group when enforcing fairness constraints.

Post-processing techniques aim to enhance fairness by adjusting the model's output distribution to match the training distribution or a specific fairness metric. This calibration is done by taking the model's prediction and the sensitive attribute as inputs \cite{hardt2016equality,zhao2017men,du2020fairness}. Although effective in reducing bias, post-processing methods require access to the sensitive attribute during inference, which may not be feasible for disease diagnosis using image inputs.

In this paper, we propose a general multi-exit (ME) training framework to enhance model fairness without the need to modify data preprocessing or inference strategy, moreover, since the early exit strategy has high confidence in its early exit instances, it does not sacrifice the accuracy of any group and maintains high overall accuracy and fairness simultaneously.

% Bias mitigation methods are designed to reduce the native bias in the dataset to reduce the chance of unfair prediction. There are largely three avenues for current debiased strategy, including pre-processing, in-processing, and post-processing. 

% Pre-processing strategies usually remove the information which may cause ``discrimination'' from training data before training. \cite{kamiran2012data} use different weight to neutralize the effect of the sensitive information in training phase and present the experiments result on real-life data. \cite{ngxande2020bias,lu2020gender} achieve fairness via data pre-process methods including data generation and data augmentation. These strategies are not suitable for dermatological data because the sensitive information which can help diagnosis (i.e., diseased area) in image may be altered or removed.  

% In-processing method usually modify on-the-shelf model architecture, loss function, and model regularization to achieve fairness goal. For example, adversarial training \cite{alvi2018turning,zhang2018mitigating} mitigates the bias through adversarially trains an encoder and classifier to learn a fair representation. Some regularization-based methods, such as \cite{quadrianto2019discovering,jung2021fair} used Hilbert-Schmidt norm and Maximum Mean Discrepancy (MMD) \cite{gretton2012kernel} respectively to learn a fair representation that retain the features' semantics from input domain. Recently, pruning techniques are proved to be an effective method to achieve fairness in dermatological disease diagnosis. \cite{wu2022fairprune} achieve fairness via pruning the parameters that hold large discrepancy of saliencies between different sensitive attributes.

% Post-processing method \cite{hardt2016equality,zhao2017men,du2020fairness} aims to calibrate the model's output to enhance fairness. They need to use sensitive attribute and prediction distribution to modify the previous distribution result. However, these methods need to use sensitive information to achieve fairness, while the disease diagnosis taking images as input.

% Compared with the above methods, we proposed a general bias mitigation framework that could leverage multiple fairness constraints


% Disentangled representation forces the encoder abandon the information of sensitive attributes and thus achieve fairness
%% ?????
% However, those pre-processing and post-processing methods are not suitable for medical image diagnosis task because they need to remove or access sensitive attribute during experiment, it may influence the diagnosis, so our method focus on improving fairness in training phase. Moreover, out method can achieve fairness without knowing sensitive attribute in the training time, which is a "fairness through unawareness" technique.

% However, most pre-processing and post-processing methods are unsuitable for medical image diagnosis tasks because they are required to remove or access sensitive attributes during the experiment, which are not practical. As a result, our method focus on improving fairness in the training phase. Moreover, our method can achieve fairness without knowing sensitive attributes during training, which is a "fairness through unawareness" technique.